\clearpage\begingroup\renewcommand{\otThumb}{}\renewcommand{\otHeadTag}{D}\phantomsection\label{dossiers}\pdfbookmark[0]{Dossiers et évaluations}{dossiers}\fiche{Une spirale et son centre}{Synthèse · Complexes et suites}\phantomsection\label{projet-1}\pdfbookmark[1]{Une spirale et son centre}{bm-projet-1}
$z_0=0$, $z_{n+1}=az_n+1$ avec $a=(1+\mathrm i)/2$.
\repere{projet-cor-1}{Solution}
\begin{questions}
\item Trouver le point fixe $\ell$.
\item Exprimer $z_n-\ell$ et sa distance à zéro.
\item Décrire l’évolution géométrique autour de $\ell$.
\item Trouver un rang garantissant $|z_n-\ell|<0{,}01$.
\end{questions}
\vspace{5mm}\souscours{Pour communiquer}\noindent Présenter les hypothèses, la démarche et une conclusion ; accompagner les résultats numériques de leur précision. Garder une trace de recherche, y compris des essais qui ont permis de choisir la méthode.\par\vfill\noindent{\otMedium\fontsize{8}{10}\selectfont\color{otMuted}Espace de recherche}\par\vspace{2mm}\noindent\begin{tikzpicture}\draw[otRule,line width=.2pt,step=5mm] (0,0) grid (160mm,60mm);\end{tikzpicture}
\fiche{Le pentagone régulier}{Synthèse · Racines de l’unité}\phantomsection\label{projet-2}\pdfbookmark[1]{Le pentagone régulier}{bm-projet-2}
On pose $\omega=\mathrm e^{2\mathrm i\pi/5}$ et $c=\cos(2\pi/5)$.
\repere{projet-cor-2}{Solution}
\begin{questions}
\item Justifier $1+\omega+\cdots+\omega^4=0$.
\item Prendre les parties réelles et obtenir $1+2c+2\cos(4\pi/5)=0$.
\item En déduire la valeur exacte de $c$.
\item Calculer la longueur d’un côté du pentagone inscrit dans le cercle unité.
\end{questions}
\vspace{5mm}\souscours{Pour communiquer}\noindent Présenter les hypothèses, la démarche et une conclusion ; accompagner les résultats numériques de leur précision. Garder une trace de recherche, y compris des essais qui ont permis de choisir la méthode.\par\vfill\noindent{\otMedium\fontsize{8}{10}\selectfont\color{otMuted}Espace de recherche}\par\vspace{2mm}\noindent\begin{tikzpicture}\draw[otRule,line width=.2pt,step=5mm] (0,0) grid (160mm,60mm);\end{tikzpicture}
\fiche{Un message affine}{Synthèse · Congruences et chiffrement}\phantomsection\label{projet-3}\pdfbookmark[1]{Un message affine}{bm-projet-3}
Les lettres sont numérotées de $0$ à $25$. On chiffre par $y\equiv7x+3\pmod{26}$.
\repere{projet-cor-3}{Solution}
\begin{questions}
\item Pourquoi le déchiffrement est-il possible ?
\item Trouver la formule inverse.
\item Chiffrer MAT.
\item Déchiffrer le résultat et discuter l’intérêt pédagogique du modèle.
\end{questions}
\vspace{5mm}\souscours{Pour communiquer}\noindent Présenter les hypothèses, la démarche et une conclusion ; accompagner les résultats numériques de leur précision. Garder une trace de recherche, y compris des essais qui ont permis de choisir la méthode.\par\vfill\noindent{\otMedium\fontsize{8}{10}\selectfont\color{otMuted}Espace de recherche}\par\vspace{2mm}\noindent\begin{tikzpicture}\draw[otRule,line width=.2pt,step=5mm] (0,0) grid (160mm,60mm);\end{tikzpicture}
\fiche{Pourquoi un RSA miniature fonctionne}{Synthèse · Fermat, Gauss et puissances}\phantomsection\label{projet-4}\pdfbookmark[1]{Pourquoi un RSA miniature fonctionne}{bm-projet-4}
On prend $p=7$, $q=11$, $n=77$, $e=7$ et $d=43$. On considère des messages entiers premiers avec $77$.
\repere{projet-cor-4}{Solution}
\begin{questions}
\item Verifier $ed\equiv1\pmod{60}$.
\item Chiffrer le message $m=2$ par $c\equiv m^e\pmod{77}$.
\item Montrer $m^{ed}\equiv m$ modulo $7$ puis modulo $11$.
\item Conclure modulo $77$ et contrôler le déchiffrement avec Python.
\end{questions}
\vspace{5mm}\souscours{Pour communiquer}\noindent Présenter les hypothèses, la démarche et une conclusion ; accompagner les résultats numériques de leur précision. Garder une trace de recherche, y compris des essais qui ont permis de choisir la méthode.\par\vfill\noindent{\otMedium\fontsize{8}{10}\selectfont\color{otMuted}Espace de recherche}\par\vspace{2mm}\noindent\begin{tikzpicture}\draw[otRule,line width=.2pt,step=5mm] (0,0) grid (160mm,60mm);\end{tikzpicture}
\fiche{Deux calendriers}{Synthèse · Équations entières}\phantomsection\label{projet-5}\pdfbookmark[1]{Deux calendriers}{bm-projet-5}
Un événement se produit aux jours $x\equiv4\pmod7$, un autre aux jours $x\equiv3\pmod5$.
\repere{projet-cor-5}{Solution}
\begin{questions}
\item Poser $x=4+7k$ et trouver $k$ modulo $5$.
\item Donner tous les jours communs.
\item Lesquels appartiennent à $[1,100]$ ?
\item Justifier la période et l’unicité de la classe obtenue.
\end{questions}
\vspace{5mm}\souscours{Pour communiquer}\noindent Présenter les hypothèses, la démarche et une conclusion ; accompagner les résultats numériques de leur précision. Garder une trace de recherche, y compris des essais qui ont permis de choisir la méthode.\par\vfill\noindent{\otMedium\fontsize{8}{10}\selectfont\color{otMuted}Espace de recherche}\par\vspace{2mm}\noindent\begin{tikzpicture}\draw[otRule,line width=.2pt,step=5mm] (0,0) grid (160mm,60mm);\end{tikzpicture}
\fiche{Visiter toutes les liaisons}{Synthèse · Graphe complet et Euler}\phantomsection\label{projet-6}\pdfbookmark[1]{Visiter toutes les liaisons}{bm-projet-6}
Cinq lieux sont reliés deux à deux. On note les sommets $0,1,2,3,4$.
\repere{projet-cor-6}{Solution}
\begin{questions}
\item Identifier le graphe, ses degrés et son nombre d’arêtes.
\item Un cycle eulérien existe-t-il ?
\item Vérifier le parcours $0,1,2,0,3,1,4,2,3,4,0$.
\item La puissance $A^{10}$ permet-elle de compter seulement les cycles eulériens ?
\end{questions}
\vspace{5mm}\souscours{Pour communiquer}\noindent Présenter les hypothèses, la démarche et une conclusion ; accompagner les résultats numériques de leur précision. Garder une trace de recherche, y compris des essais qui ont permis de choisir la méthode.\par\vfill\noindent{\otMedium\fontsize{8}{10}\selectfont\color{otMuted}Espace de recherche}\par\vspace{2mm}\noindent\begin{tikzpicture}\draw[otRule,line width=.2pt,step=5mm] (0,0) grid (160mm,60mm);\end{tikzpicture}
\fiche{Une parabole par trois mesures}{Synthèse · Matrices et interpolation}\phantomsection\label{projet-7}\pdfbookmark[1]{Une parabole par trois mesures}{bm-projet-7}
On cherche $f(x)=ax^2+bx+c$ passant par $(0,1),(1,3),(2,9)$.
\repere{projet-cor-7}{Solution}
\begin{questions}
\item Écrire le système sur $(c,b,a)$.
\item Donner sa matrice et expliquer l’unicité.
\item Résoudre et vérifier.
\item Un quatrième point non compatible pourrait-il être imposé sans changer le degré ?
\end{questions}
\vspace{5mm}\souscours{Pour communiquer}\noindent Présenter les hypothèses, la démarche et une conclusion ; accompagner les résultats numériques de leur précision. Garder une trace de recherche, y compris des essais qui ont permis de choisir la méthode.\par\vfill\noindent{\otMedium\fontsize{8}{10}\selectfont\color{otMuted}Espace de recherche}\par\vspace{2mm}\noindent\begin{tikzpicture}\draw[otRule,line width=.2pt,step=5mm] (0,0) grid (160mm,60mm);\end{tikzpicture}
\fiche{Diffusion d’Ehrenfest à deux particules}{Synthèse · Chaîne à trois états}\phantomsection\label{projet-8}\pdfbookmark[1]{Diffusion d’Ehrenfest à deux particules}{bm-projet-8}
Deux particules sont réparties entre deux boîtes. À chaque étape, on choisit uniformément une particule et on la change de boîte. L’état est le nombre de particules dans la première boîte.
\repere{projet-cor-8}{Solution}
\begin{questions}
\item Construire la matrice $R$ dans l’ordre $0,1,2$.
\item Déterminer une distribution invariante.
\item Partant de $0$, montrer l’absence de convergence.
\item On reste maintenant immobile avec probabilité $1/2$, sinon on applique une transition de $R$. Construire $P$ et comparer au cours.
\end{questions}
\vspace{5mm}\souscours{Pour communiquer}\noindent Présenter les hypothèses, la démarche et une conclusion ; accompagner les résultats numériques de leur précision. Garder une trace de recherche, y compris des essais qui ont permis de choisir la méthode.\par\vfill\noindent{\otMedium\fontsize{8}{10}\selectfont\color{otMuted}Espace de recherche}\par\vspace{2mm}\noindent\begin{tikzpicture}\draw[otRule,line width=.2pt,step=5mm] (0,0) grid (160mm,60mm);\end{tikzpicture}
\fiche{Un modèle de classement de pages}{Synthèse · Marche aléatoire et redémarrage}\phantomsection\label{projet-9}\pdfbookmark[1]{Un modèle de classement de pages}{bm-projet-9}
Dans un réseau fictif, $A$ pointe vers $B,C$ à parts égales ; $B,C$ pointent vers $A$. À chaque étape, on suit un lien avec probabilité $0{,}85$, sinon on choisit une page parmi les trois uniformément.
\repere{projet-cor-9}{Solution}
\begin{questions}
\item Construire la matrice $P$.
\item Justifier que les probabilités d’arrivée en $B,C$ sont égales après une étape.
\item Trouver une distribution invariante et classer les pages.
\item Établir une récurrence pour la probabilité $a_n$ d’être en $A$ et discuter la convergence.
\end{questions}
\vspace{5mm}\souscours{Pour communiquer}\noindent Présenter les hypothèses, la démarche et une conclusion ; accompagner les résultats numériques de leur précision. Garder une trace de recherche, y compris des essais qui ont permis de choisir la méthode.\par\vfill\noindent{\otMedium\fontsize{8}{10}\selectfont\color{otMuted}Espace de recherche}\par\vspace{2mm}\noindent\begin{tikzpicture}\draw[otRule,line width=.2pt,step=5mm] (0,0) grid (160mm,60mm);\end{tikzpicture}
\fiche{Évaluation 1 · Nombres complexes}{60 minutes · 20 points}\phantomsection\label{evaluation-1}\pdfbookmark[1]{Évaluation 1 · Nombres complexes}{bm-evaluation-1}
On considère $P(z)=z^3-z^2+4z-4$ et les points d’affixes $a=0,b=2,c=2\mathrm i$.
\repere{evaluation-cor-1}{Solution}
\begin{questions}
\item Factoriser $P$ à partir de la racine $1$ et résoudre $P(z)=0$. \hfill(4 points)
\item Donner module, argument et forme exponentielle de $1+\mathrm i$ ; calculer sa puissance $12$. \hfill(4 points)
\item Interpréter $(c-a)/(b-a)$ et décrire le triangle. \hfill(4 points)
\item Déterminer et placer les racines sixièmes de l’unité. \hfill(4 points)
\item Retrouver $\cos3x=4\cos^3x-3\cos x$ et calculer $\int_0^{\pi/2}\cos^3x\,\mathrm dx$. \hfill(4 points)
\end{questions}
\vspace{5mm}\souscours{Pour communiquer}\noindent Présenter les hypothèses, la démarche et une conclusion ; accompagner les résultats numériques de leur précision. Garder une trace de recherche, y compris des essais qui ont permis de choisir la méthode.\par\vfill\noindent{\otMedium\fontsize{8}{10}\selectfont\color{otMuted}Espace de recherche}\par\vspace{2mm}\noindent\begin{tikzpicture}\draw[otRule,line width=.2pt,step=5mm] (0,0) grid (160mm,60mm);\end{tikzpicture}
\fiche{Évaluation 2 · Arithmétique}{60 minutes · 20 points}\phantomsection\label{evaluation-2}\pdfbookmark[1]{Évaluation 2 · Arithmétique}{bm-evaluation-2}
Tous les inconnus sont entiers, sauf indication contraire.
\repere{evaluation-cor-2}{Solution}
\begin{questions}
\item Calculer le PGCD de $252,198$ et donner un couple de Bézout. \hfill(4 points)
\item Résoudre $7x\equiv5\pmod{26}$. \hfill(4 points)
\item Résoudre $15x+21y=6$ dans $\Z^2$. \hfill(4 points)
\item Étudier la primalité de $223$, puis calculer $2^{100}$ modulo $13$. \hfill(4 points)
\item Prouver l’infinitude des nombres premiers sans supposer que le produit d’une liste de premiers augmenté de $1$ est premier. \hfill(4 points)
\end{questions}
\vspace{5mm}\souscours{Pour communiquer}\noindent Présenter les hypothèses, la démarche et une conclusion ; accompagner les résultats numériques de leur précision. Garder une trace de recherche, y compris des essais qui ont permis de choisir la méthode.\par\vfill\noindent{\otMedium\fontsize{8}{10}\selectfont\color{otMuted}Espace de recherche}\par\vspace{2mm}\noindent\begin{tikzpicture}\draw[otRule,line width=.2pt,step=5mm] (0,0) grid (160mm,60mm);\end{tikzpicture}
\fiche{Évaluation 3 · Matrices et graphes}{75 minutes · 20 points}\phantomsection\label{evaluation-3}\pdfbookmark[1]{Évaluation 3 · Matrices et graphes}{bm-evaluation-3}
$A=\begin{pmatrix}1&1\\0&1\end{pmatrix}$. Une chaîne indépendante de cette première matrice a pour transition $P=\begin{pmatrix}0{,}8&0{,}2\\0{,}3&0{,}7\end{pmatrix}$.
\repere{evaluation-cor-3}{Solution}
\begin{questions}
\item Donner et prouver $A^n$ par récurrence. \hfill(4 points)
\item Résoudre $AX=\binom53$ puis calculer $A^{10}\binom01$. \hfill(4 points)
\item Construire le graphe de transition de $P$ et calculer $P^2$. \hfill(4 points)
\item Depuis $\pi_0=(1;0)$, déterminer la distribution après $n$ transitions et sa limite. \hfill(4 points)
\item Expliquer la différence entre le coefficient d’une puissance de matrice d’adjacence et celui d’une puissance de matrice de transition. \hfill(4 points)
\end{questions}
\vspace{5mm}\souscours{Pour communiquer}\noindent Présenter les hypothèses, la démarche et une conclusion ; accompagner les résultats numériques de leur précision. Garder une trace de recherche, y compris des essais qui ont permis de choisir la méthode.\par\vfill\noindent{\otMedium\fontsize{8}{10}\selectfont\color{otMuted}Espace de recherche}\par\vspace{2mm}\noindent\begin{tikzpicture}\draw[otRule,line width=.2pt,step=5mm] (0,0) grid (160mm,60mm);\end{tikzpicture}
\clearpage\endgroup